\section{Limitations}
\label{sec:limitations}

\paragraph{Private pseudo-reference.}
The private triangulated evaluator uses 2D evidence from the same videos as the
monocular 3D inputs. Its errors can be correlated with candidate errors, and
sequence-level Sim3 plus hip centering removes some failure modes before
scoring. Private MPJPE is agreement with one upstream construction, not
absolute 3D accuracy.

\paragraph{Small single-seed held-out study.}
The primary learned test contains only 14 people and one training seed. Paired
tests control person reuse across methods but do not estimate optimization
variance or broad population uncertainty. The 137-person learned values include
training and validation participants and are descriptive only.

\paragraph{Limited synthetic external validation.}
Unity provides native 3D independent of the private pseudo-reference, but only
199 samples, three sequences, one avatar/environment, and one fixed camera
pair. It cannot establish population generality. Its negative A6 result and the
advantage of calibrated triangulation should be replicated on diverse real
multi-view benchmarks.

\paragraph{Observational cohort labels and unmeasured confounding.}
Elderly and student are dataset group labels, not randomized exposures.
Individual age, sex, health, training history, body size, and other covariates
are unavailable. Person-disjoint inference prevents direct train/test leakage
but does not make the comparison causal. The reported coefficients cannot be interpreted as ageing effects,
normative values, impairment markers, or clinical evidence.

\paragraph{Pose-source and estimand sensitivity.}
The A6 speed and jerk coefficients are much larger than the same centered
mixed-model coefficients from face, side, or deterministic features.
Person-median A6 comparisons are also not corrected-significant. This
instability limits the application to hypothesis generation until independent
kinematics and repeated cross-fit seeds are available.

\paragraph{Dependence on upstream errors.}
Both views originate from one monocular estimator and skeletal model.
Canonicalization cannot correct shared joint-definition, anthropometric, or
hallucination errors. High-consensus identity may reinforce common bias, and
synthetic corruptions cannot span the full real error distribution.

\paragraph{Alignment, views, and movement specificity.}
The loader honors one global side-to-face offset and does not model drift,
dropped frames, rolling shutter, or nonlinear local misalignment. Only one
face/side pair and one gymnastics movement family are studied. The architecture
is symmetric to view order but does not support a variable number of cameras.

\paragraph{Coordinate and anatomical interpretation.}
The output is restored to an uncalibrated face reference. Pelvis and thorax
frames depend directly on estimated hips and shoulders, and A3-relative
retention is not a comparison to physical motion. The method does not estimate
validated clinical joint angles, forces, or moments.

\paragraph{Model and governance choices.}
Loss weights, residual bounds, and corruption probabilities are manually
specified. Ethics approval or exemption, participant consent wording,
institutional postal address, and release conditions for code and derived data
remain author-governance items that must be verified before submission.
